\documentclass[../../../include/open-logic-section]{subfiles}

\begin{document}

\olfileid{sth}{spine}{recursion}
\olsection{超限递归定理}

然而，我们必须首先确认 \olref[valpha]{defValphas} 是一个成功的定义。更一般地，我们需要证明，任何通过（超限）递归给出超限定义的尝试都会成功。这正是本节的目标。

\emph{警告：这部分内容比较复杂}。不过，其核心要旨非常简单：超限归纳加上替换公理保证了（若干版本的）超限递归的合法性。\footnote{提醒：所有公式和项都可以带有参数（除非另有明确说明）。}

\begin{defn}
	设 $\tau(x)$ 是一个项；设 $f$ 是一个函数；设 $\alpha$ 是一个序数。我们称 $f$ 是 $\tau$ 的一个 \emph{$\alpha$-近似}，当且仅当 $\dom{f} = \alpha$ 且 $(\forall \beta \in \alpha)f(\beta) = \tau(\funrestrictionto{f}{\beta})$。
\end{defn}

\begin{lem}[有界递归]\ollabel{transrecursionfun}
对任意项 $\tau(x)$ 和任意序数 $\alpha$，存在唯一的 $\tau$ 的 $\alpha$-近似。
\end{lem}
\begin{proof}
我们将证明，对任意 $\gamma \leq \alpha$，都存在唯一的 $\gamma$-近似。

我们首先证明唯一性。设 $g$ 和 $h$（分别）是 $\gamma$-近似和 $\delta$-近似。对其自变元进行超限归纳可知，对任意 $\beta \in \dom{g} \cap \dom{h} =
\gamma \cap \delta = \min(\gamma, \delta)$ 都有 $g(\beta) = h(\beta)$。因此我们的近似（如果存在的话）是唯一的，并且在所有值上都是一致的。

为证明存在性，我们现在对序数 $\delta \leq \alpha$ 施加简单的超限归纳（\olref[ordinals][opps]{simpletransrecursion}）。

空函数显然是一个 $\emptyset$-近似。

如果 $g$ 是一个 $\gamma$-近似，那么 $g \cup
\{\tuple{\gamma, \tau(g)}\}$ 是一个 $\ordsucc{\gamma}$-近似。

如果 $\gamma$ 是一个极限序数，并且对所有的 $\delta < \gamma$，$g_\delta$ 都是 $\delta$-近似，则令 $g = \bigcup_{\delta \in \gamma} g_\delta$。由于各个 $g_\delta$ 在所有值上都是一致的，所以这是一个函数。并且如果 $\delta \in \gamma$，那么 $g(\delta) = g_{\ordsucc{\delta}}(\delta) = \tau(\funrestrictionto{g_{\ordsucc{\delta}}}{\delta}) = \tau(\funrestrictionto{g}{\delta})$。

这就完成了超限归纳证明。
\end{proof}

如果我们定义的是一个 \emph{项} 而非函数，那么我们可以从上述结果中去除界限 $\alpha$。在下面结果的陈述和证明中，当 $\sigma$ 是一个项时，我们令 $\funrestrictionto{\sigma}{\alpha} = \Setabs{\tuple{\beta,
\sigma(\beta)}}{\beta \in \alpha}$。

\begin{thm}[一般递归]\ollabel{transrecursionschema}
对任意项 $\tau(x)$，我们可以显式地定义一个项 $\sigma(x)$，使得对任意序数~$\alpha$，都有 $\sigma(\alpha) = \tau(\funrestrictionto{\sigma}{\alpha})$。
\end{thm}

\begin{proof}
对每个 $\alpha$，根据 \olref{transrecursionfun}，存在唯一的 $\tau$ 的 $\alpha$-近似 $f_\alpha$。定义 $\sigma(\alpha)$ 为 $f_{\ordsucc{\alpha}}(\alpha)$。于是：
	\begin{align*}
		\sigma(\alpha) &=
		f_{\ordsucc{\alpha}}(\alpha) \\&=
		\tau(\funrestrictionto{f_{\ordsucc{\alpha}}}{\alpha}) \\&=
		\tau(\Setabs{\tuple{\beta, f_{\ordsucc{\alpha}}(\beta)}}{\beta \in \alpha}) \\&=
		\tau(\Setabs{\tuple{\beta, f_{\ordsucc{\beta}}(\beta)}}{\beta \in \alpha})\\&=
		\tau(\funrestrictionto{\sigma}{\alpha})
	\end{align*}
注意到对所有的 $\beta < \alpha$，$f_{\ordsucc{\beta}}(\beta) = f_{\ordsucc{\alpha}}(\beta)$，如 \olref{transrecursionfun} 中所述。
\end{proof}
\noindent
注意，\olref{transrecursionschema} 是一个\emph{模式}。关键是，我们不能期望 $\sigma$ 定义一个函数，即某种\emph{集合}，因为那样的话 $\dom{\sigma}$ 就会是所有序数的集合，这与 Burali-Forti 悖论（\olref[ordinals][basic]{buraliforti}）矛盾。

不过，我们仍然需要证明，\olref{transrecursionschema} 能够证实我们对 $V_\alpha$ 的定义。这可能并非显而易见；但通过超限递归的最后一种简单版本，这一点将变得显然。

\begin{thm}[简单递归]\ollabel{simplerecursionschema}
对任意项 $\tau(x)$ 和 $\theta(x)$ 以及任意集合 $A$，我们可以显式地定义一个项 $\sigma(x)$，使得：
\begin{align*}
	\sigma(\emptyset) &= A\\
	\sigma(\ordsucc{\alpha}) &= \tau(\sigma(\alpha)) &&
		\text{对任意序数 }\alpha\\
	\sigma(\alpha) &= \theta(\ran{\funrestrictionto{\sigma}{\alpha}})&&
	\text{当 }\alpha\text{ 是极限序数}
\end{align*}
\end{thm}

\begin{proof}
我们首先定义一个项 $\xi(x)$ 如下：
%t $\xi(x) = A$ if $x$ is not a function whose domain is an ordinal;
%otherwise let:
\[
	\xi(x) =
	\begin{cases}
			A &\text{若 $x$ 不是}\\
			&\hspace{1em}\text{定义域为某个序数的函数；否则：}\\
			\tau(x(\alpha)) & \text{若 $\dom{x} = \ordsucc{\alpha}$}\\
			\theta(\ran{x}) & \text{若 $\dom{x}$ 是极限序数}
	\end{cases}
\]
根据 \olref{transrecursionschema}，存在一个项 $\sigma(x)$，使得对每个序数 $\alpha$，$\sigma(\alpha) = \xi(\funrestrictionto{\sigma}{\alpha})$；并且，$\funrestrictionto{\sigma}{\alpha}$ 是一个定义域为 $\alpha$ 的函数。我们通过简单的超限归纳（\olref[ordinals][opps]{simpletransrecursion}）证明 $\sigma$ 具有所需的性质。

首先，$\sigma(\emptyset) = \xi(\emptyset) = A$。

其次，$\sigma(\ordsucc{\alpha}) = \xi(\funrestrictionto{\sigma}{\ordsucc{\alpha}}) = \tau(\funrestrictionto{\sigma}{\ordsucc{\alpha}}(\alpha)) = \tau(\sigma(\alpha))$。

最后，当 $\alpha$ 是极限序数时，$\sigma(\alpha) = \xi(\funrestrictionto{\sigma}{\alpha}) = \theta(\ran{\funrestrictionto{\sigma}{\alpha}})$。
\end{proof}
\noindent
现在，为了证实 \olref[valpha]{defValphas}，只需令 $A = \emptyset$，$\tau(x) = \Pow{x}$，$\theta(x) = \bigcup x$。终于，这证实了 $V_\alpha$ 的定义！{}

\end{document}